\begin{abstract}
Deploying vision-based reinforcement learning policies in safety-critical domains requires anticipating performance collapse under visual distribution shifts without executing costly, potentially catastrophic test-time rollouts.
In this work, we rigorously evaluate representation-distance forecasting and test-time adaptation across three targeted empirical phases on a visual continuous control task (Visual Pursuit, continuous target tracking at $64\times 64$).
First, we investigate classical representation drift $d_\mathrm{rep}$ and evaluate whether semantic spatial weighting $\tilde{d}_\mathrm{rep}$ resolves metric collapse under complex visual shifts; we observe that semantic weighting provides only modest improvement (best overall Spearman $\rho = 0.1075$, $p = 0.0315$ with ground-truth segmentation masks, compared to baseline $\rho = 0.0926$, $p = 0.0643$), indicating that spatial masking alone is insufficient for robust forecasting on non-smooth variations.
Second, Riemannian analysis on the Grassmann manifold $\mathrm{Gr}(k, 512)$ reveals that end-to-end trained CNN encoders maintain a more compact representation subspace under shift, yielding significantly lower geodesic drift ($1.5402$--$2.2297$\,rad across $k \in \{5, 10, 20\}$) compared to frozen ResNet-18 backbones ($2.9735$--$5.7208$\,rad).
Third, we evaluate Test-Time Self-Supervised Adaptation (TTSS) combining temporal smoothness and moment matching with a trust-region safeguard under severe perceptual shift. TTSS reduces catastrophic tracking floor-collapse from $40.0\%$ to $20.0\%$ (a $50\%$ relative reduction). However, we find that the trust-region safeguard accepted $100.0\%$ of proposed gradient updates, demonstrating that the empirical benefit stems from the regularized self-supervised objective rather than active step rejection.
We report our findings transparently, highlighting key limitations and outlining requirements for generalizable pre-rollout competence boundary forecasting.
\end{abstract}
